\section{Existing Range Filters}
\label{sec:existing-work}

We survey the five practical range filters most directly comparable to our work.
We describe each by the technique that determines its
space--accuracy trade-off and contrast it with the Goswami lower
bound (Section~\ref{sec:bounds}).

\paragraph{Trie-based: SuRF.}
The Succinct Range Filter~\citep{zhang2018surf} stores the set of keys
as a Fast Succinct Trie --- a compact bitmap-based succinct trie
encoding at about $10$ bits per node~\citep[\S4.1.2]{zhang2018surf}.
A range query descends the trie to the boundary of the query interval.
The cost of the query is set by the trie depth, which reflects the
key length in the universe (e.g.\ $\leq 8$ levels for $64$-bit
integers under a byte-wide fanout), but \emph{not} the query length
$\mathcal{L}$; and the FPR has no explicit dependence on $\mathcal{L}$
either --- unlike a linear-probe Bloom filter, which would give
$\mathrm{FPR} \approx \mathcal{L}\varepsilon$ for a range of length $\mathcal{L}$.
This makes SuRF suitable for queries of almost any
length without a per-query penalty.

The price of this flexibility is the absence of a worst-case FPR guarantee.
A false positive occurs when the prefix of a query key
coincides with a stored key prefix, so the FPR is set by the trie
depth and by the \emph{correlation between queries and data}. The
authors themselves acknowledge that ``the point query FPR of
SuRF-Base [can be] $100\%$'' in the adversarial
case~\citep[\S3.2]{zhang2018surf}; in practice queries and keys are
``almost always correlated to the stored
keys''~\citep[\S3.1]{zhang2018surf}, and the empirical FPR on the YCSB
workload is $4\%$ on $64$-bit integer keys and $25\%$ on email
addresses~\citep[\S4.2]{zhang2018surf}.
The space is not parameterized by $\varepsilon$ either: SuRF-Base costs about $10$ bits per key for
$64$-bit integers and $14$ bits per key for emails regardless of the
target FPR.

The headline RocksDB result --- up to $5\times$ throughput on
closed-seek queries with SuRF-Real over the Bloom-filter
configuration, on a synthetic $100$~GB time-series workload with
$128$-bit timestamp$+$sensor keys, all filters sized to $14$~bits per
key~\citep[\S6, Fig.~13]{zhang2018surf} --- is an \emph{end-to-end}
effect of skipped SSTable I/O, not a faster filter operation.
A Bloom filter cannot prune range queries at all (``the Bloom filter
does not help range queries and is equivalent to having no
filter''~\citep[\S6]{zhang2018surf}),
so without a range filter RocksDB fetches candidate blocks from every
level; SuRF eliminates this I/O. The $5\times$ holds when $99\%$ of
queries return empty; at $10\%$ empty the gap narrows to roughly
$1.5\times$ (read from Fig.~13).

SuRF does not position itself against an information-theoretic
range-filter space bound and does not cite the Goswami construction.

\paragraph{Bloom hierarchy: Rosetta.}
Rosetta~\citep{luo2020rosetta} indexes every binary prefix of every
key in a Bloom hierarchy --- one filter per prefix length,
$\log_2 \mathcal{L} + 1$ filters in total when a query-length bound
$\mathcal{L}$ is known and as many as the key length
otherwise~\citep[\S2.1, \S3.1]{luo2020rosetta}. A range query of
length $\mathcal{L}$ decomposes into at most $2 \log_2 \mathcal{L}$
dyadic intervals; for each one, Rosetta probes the Bloom filter at
the root of the corresponding sub-tree and, on a positive answer,
recursively descends into deeper
levels~\citep[\S2.2.1]{luo2020rosetta}. Because Bloom filters are
insensitive to the key distribution, the FPR guarantee holds under
any workload.
Under the ``first-cut'' bit-allocation
policy~\citep[\S2.3]{luo2020rosetta}, the hierarchy occupies
$1.44 \cdot n \log_2(\mathcal{L}/\varepsilon)$ bits --- the Goswami
lower bound multiplied by the standard $1/\ln 2$ Bloom-filter
overhead~\citep[\S3.1]{luo2020rosetta}; the deployed system
redistributes the per-level budget using runtime query statistics
gathered between LSM compactions~\citep[\S2.4]{luo2020rosetta}.
The authors are explicit about the cost: ``due to the multiple
probes required (for every prefix) Rosetta has a high probe
cost''~\citep[\S1]{luo2020rosetta}. On a $50$M uint64-key benchmark
at $22$ bits per key, Rosetta is up to $4\times$ faster end-to-end
than SuRF and up to $40\times$ faster than default RocksDB on short
and medium range queries~\citep[\S5, Fig.~5(A1), (D)]{luo2020rosetta};
the advantage narrows on long ranges.

\paragraph{Learned: SNARF.}
SNARF~\citep{vaidya2022snarf} fits a piecewise-linear monotonic
model $F$ to the empirical Cumulative Distribution Function (CDF) of
the key set and uses it to map each stored key $x$ to bit position
$\lfloor F(x) \cdot nK \rfloor$ in a sparse bit array of size $nK$.
The set bit positions are then stored either with
Golomb coding~\citep{golomb1966runlength} (default, close to the
information-theoretic optimum) or with Elias--Fano coding (faster queries at slightly
higher space cost).
The compressed array is split into fixed-size
segments so that a query only decompresses the segments overlapping
its interval~\citep[\S2.3.2]{vaidya2022snarf}.
The parameter $K$ controls the bit array size ($|B| = nK$); total
space is $\approx 2.4 + \log_2 K$ bits per
key~\citep[\S3]{vaidya2022snarf}.

On already-sorted input SNARF builds in $O(n)$: sampling, fitting
linear segments between adjacent sample points, filling the bit
array and compressing it are all linear.
A range query costs $O(\log n + S)$, where $\log n$
is the binary search in the model's sample array and $S$ is the
number of decompressed segments.

SNARF's FPR analysis rests on two assumptions: the model maps
keys approximately uniformly into the bit array, and queries are
uncorrelated with stored keys. Under these assumptions FPR
$\approx 1/K$ for point queries and range queries, with no
explicit dependence on the range
length~\citep[\S3]{vaidya2022snarf}. Setting $K \approx
1/\varepsilon$ gives total space $\approx 2.4 +
\log_2(1/\varepsilon)$ bits per key, which absorbs the
$\log_2 \mathcal{L}$ term of the Goswami bound: SNARF goes below the
bound when its assumptions about the data hold.
On the Search on Sorted Data (SOSD)
benchmark~\citep{kipf2019sosd} SNARF reports up to $50\times$
better FPR than SuRF and $100\times$ better FPR than Rosetta at
matched memory~\citep[\S2]{vaidya2022snarf}.

Both assumptions break in practice.
The paper itself warns about correlated queries: ``when a query end point is consistently
close to an actual key but the query does not include a key, it
may yield a false positive in SNARF and
SuRF''~\citep[\S4.1]{vaidya2022snarf}, and recommends Rosetta for
short, highly correlated range queries. The uniform mapping
assumption fails empirically on real-world distributions with
structured sparsity. On the OpenStreetMap (OSM) dataset from SOSD
--- cell IDs spanning nearly the full 64-bit universe yet
concentrated in a small fraction of it --- the paper reports that
SNARF reaches FPR $\leq 10^{-4}$ at $< 10$ bits per key only
``for most cases'', and Fig.~5(B) shows the OSM range-query curve
flooring well above the $10^{-4}$ level reached on the other SOSD
datasets, even at $16$ bits per
key~\citep[\S6.1.2, Fig.~5(B)]{vaidya2022snarf}.

\paragraph{Practical realisation of Goswami: Grafite.}
Grafite~\citep{grafite2024} is the closest related work to this thesis.
The authors explicitly frame it as a simplified, practical
instantiation of the Goswami construction~\citep[\S3]{grafite2024}:
$h(x) = (q(\lfloor x/r \rfloor) + x) \bmod r$ with
pairwise-independent $q$ maps the universe $U$ down to a smaller
range $[0, r)$, with $r = n\mathcal{L}/\varepsilon$; the resulting hash
codes are stored using Elias--Fano coding, and a range query reduces
to a single predecessor probe~\citep[Algorithm 2]{grafite2024}.
Theorem 3.4 establishes a space bound of
$n \log_2(\mathcal{L}/\varepsilon) + 2n + o(n)$ bits and a query
time of $O(\log \mathcal{L}/\varepsilon)$ --- matching the Goswami
lower bound to within $O(n)$ bits.
Construction takes $O(n \log n)$ time, dominated by sorting the
hash codes prior to the Elias--Fano
encoding~\citep[Algorithm 1]{grafite2024}.

The authors' own framing of the headline result reads:
``Grafite is the only design to date to achieve a robust and predictable false
positive rate across all combinations of datasets, query workloads,
and range sizes, while also providing faster queries and
construction times, and dominating all competitors in the case of
correlated query workloads''~\citep[\S1]{grafite2024}.

Because Grafite is the main baseline of this thesis, we devote a
separate detailed analysis to it in
Section~\ref{sec:grafite-deep-dive} --- internals, parameter
choices, and improvements we identified during the empirical
comparison.

\paragraph{Dynamic extension: Memento.}
Memento~\citep{eslami2024memento} is the SIGMOD~2024 filter that is
the first to add dynamicity --- insertions, deletions, and growth
--- to range filtering while preserving a robust FPR guarantee for
any workload~\citep[\S1, \S4]{eslami2024memento}.
It is built on top of a Rank-and-Select Quotient Filter: each \emph{keepsake
box} --- a group of keys with a common fingerprint within one run
--- stores that fingerprint together with a contiguous list of
fixed-length key suffixes (``mementos''), all packed via Robin Hood
hashing~\citep[\S3, \S4]{eslami2024memento}.

On static workloads Memento's FPR is up to $1.5\times$ higher than
Grafite's at matched memory~\citep[\S7.1, Fig.~9]{eslami2024memento}.
The small gap reflects Memento's ${\sim}3.3n$-bit metadata overhead
versus Grafite's $2n$~\citep[\S6]{eslami2024memento}.
Range queries take $O(\log_2 \mathcal{L})$ CPU operations and
$1$--$2$ random memory accesses, both holding against any
workload~\citep[\S6, Table~2]{eslami2024memento}. This matches
Grafite on the cache-miss bound; in practice Memento is slightly
slower than Grafite~\citep[\S7.1]{eslami2024memento}.

The authors explicitly note that static range filters suffice for the immutable files of an
LSM-Tree deployment~\citep[\S2]{eslami2024memento}; Memento targets
workloads where the key set itself is mutable --- single-global-filter
LSM-trees, B-Tree indexes, and systems that mix transactional and analytical workloads.
Because our thesis stays within the static setting, Memento is
orthogonal rather than a head-to-head competitor; on the static axis
it lands close to Grafite.
